\section{Conclusion}
\label{sec:conclusion}

\vfour{} was published on the strength of a panel of weaker opponents. Played
against a copy of itself it spends \numDeadAsk\% of its asks on cards it can
\emph{prove} the target does not hold, repeats \numRepeatAsk\% of its asks
exactly, and gets \numDeclWrong\% of its declarations wrong. This report
measures that failure, corrects the account the v0.4 study gave of it, and
reports the policy that removes it. What follows is what the evidence in
\S\ref{sec:diagnosis} and \S\ref{sec:results} supports, stated without
extension.

\paragraph{The deadlock is a fixed point of the ask policy, not a freeze of
information.} In a game where nothing compels a claim, the natural explanation
for two patient policies stalling is that they have run out of things to learn.
That explanation is wrong here twice over. It is wrong about the mechanism: the
observed stall is a deterministic two-question cycle, present in \numPureCycle{}
of \numLongGames{} long games at the forensic seed, and across the pooled seed
banks no half-suit is locked to any team at any point in the cycle in
\numNoLockPct\% of long games. And it is wrong about
information: where a half-suit \emph{is} locked, \numCertRate\% of the asks its
owning team may legally make inside it strictly raise a teammate's exact
probability of the correct allocation, by a mean of \numCertMean{} and up to
\numCertMax{}. Ownership monotonicity does not imply informational monotonicity,
and the two were conflated.

The practical consequence is that the repair is a subtraction, not a forcing
rule. Declining asks whose success probability is exactly zero --- a hard
deduction from the actor's own public knowledge, which cannot introduce a
modelling error --- takes the dead-ask rate from \numFourDeadAsk\% to
\numFiveDeadAsk\%, the share of games containing a run of six consecutive dead
asks from \numFourDeadRunGames\% to \numFiveDeadRunGames\%, and the longest such
run from \numFourDeadRunLongest{} to \numFiveDeadRunLongest{}. It does so with
the event-count escalation that used to force claims switched \emph{off}:
declarations made at or after the old horizon fall from
\numFourPostHorizonDecl{} to \numFivePostHorizonDecl{}, declaration error from
\numFourDeclWrong\% to \numFiveDeclWrong\%, and no game in either arm reaches
the action cap. Forcing claims was never the fix; it was the mechanism that
converted the stall into misdeclaration, at a measured cost of
\numStageTwoCost{} points against a mirror opponent.

\paragraph{A per-card argmax allocation is infeasible, not merely unlikely.}
The second defect is of a different kind from a low-probability guess, and the
difference is worth stating exactly. \vfour{} named its forced-endgame
allocation by taking, card by card, the teammate with the highest marginal ---
a procedure with no capacity constraint, which can therefore give a teammate
more cards than that teammate holds. All \numForcedN{} such declarations over
\numForcedGames{} distinct mirror games were wrong, and the reason is not that
the estimate was poor. The named allocation has posterior probability exactly
zero: it is excluded by the declarer's own arithmetic, not merely disfavoured by
it. A policy losing a hard inference problem and a policy naming an outcome it
has itself ruled out are different failures and need different repairs. Searching
the capacity-feasible set instead takes forced-endgame accuracy from
\numVfourForcedAcc\% to \numVfiveForcedAcc\% over \numForcedEightGames{} games
per arm --- against a measured ceiling of \numFeasibleRight\% for the best
feasible allocation at those states, so the gap is narrowed and not closed. The
state is also reached \numForcedRateRatio$\times$ less often, but that belongs to
M1 rather than to M2: what was stripping a team of cards while half-suits stood
undeclared was the runs of guaranteed misses, not the declaration rule
(\S\ref{sec:results-forced}).

\paragraph{What the v0.4 study said that is withdrawn.}
\S\ref{sec:corrections} is the register. The load-bearing entry is C1, the
information-freeze account of non-termination, which is retracted with the
measurement that refutes it; the v0.4 \emph{paper} states the correct position
and is not implicated, while the supporting artifact states the wrong one. The
remaining entries correct the reported termination property (purchased, and the
price unreported), the status of the forced-endgame accuracy figure (a
mechanical certainty, not a hard problem lost), the sign of the value-function
``known gap,'' the claim that the policy prior is two channels, the usability of
asks inside a locked half-suit, the willingness ladder, the well-posedness of the
stopping problem, and the scope of the published worst case. Three of those
entries also carry corrections to earlier drafts of the present study rather
than to the v0.4 one, and are recorded in the same place for the same reason.

\paragraph{What \vfive{} costs, and what it does not buy.} It does not buy win
rate. Over the nine-style panel of \S\ref{sec:results}
(\numStyleDeals{} deals in six rotations per cell, seed \numStyleSeed{}) the
means are level --- \numVfiveMean\% for \vfive{} against \numVfourMean\% for
\vfour{} --- and head to head \vfive{} scores \numPanelVfiveVfour\%
[\numStyleCIVfour], which is $\numDeltaVfour$ points on \vfour{}'s
\numPanelVfourVfour\% mirror baseline on an interval that contains even money.
Across \numHeadSeeds{} further seed banks the same pairing averages \numHeadWin\%
with a range of \numHeadRange\%. On minimax regret over the style set \vfour{} is
marginally the better of the two, \numVfourRegret{} against \vfive{}'s
\numVfiveRegret{}. Per style the differences run in both directions:
$\numDeltaLockout$ on the lockout opponent, $\numDeltaVfour$ on the mirror and
$\numDeltaDiversifier$ on the diversifier against $\numDeltaVtwo$ on \vtwo{} and
$\numDeltaVthree$ on \vthree{}, with the remainder at half a point or less and no
per-cell interval separating either policy from the other on any style.
\vfive{}'s worst case over the nine styles is \numVfiveWorst\%, against \vfour{},
and \vfour{}'s is \numVfourWorst\%, against a copy of itself. In both cases the
worst cell is the opponent of the policy's own strength, which is the worst case
for any policy in this family, and which is why the v0.4 study's floor of
\numVfourVsVthree\% describes a panel rather than a policy. The case for \vfive{} is the pathology table and the
deception panel below. It is not the win rate, and this report does not offer
one.

\paragraph{The result that speaks to the incident that started this.} The work
was prompted by a live report: the project owner held cards of a half-suit he
had been asked for and then declined to ask back in it, and \vfour{} was visibly
confused. Modelled as an archetype --- after being asked in half-suit $S$ while
holding other cards of $S$, avoid asking in $S$ for the next several turns ---
\vfour{} wins \numDecWithholderFour\% of its games against that archetype and
\vfive{} wins \numDecWithholderFive\%, a gain of $\numDecWithholderDelta$ points
replicated at \numDecSeeds{} independent seed banks of \numDecGames{} games
each. The silent archetype moves the same way, by $\numDecSilentDelta$ points.

The mechanism is worth being precise about, because it is not the one the design
document proposed. \vfive{} has no opponent model that \vfour{} lacks: M3--M7
are specified and unbuilt, the policy prior is still fixed and global rather than
per-seat, and nothing in \vfive{} watches a seat and updates a type. What changed is that
\vfive{} never spends a turn on an ask it can prove will miss. A misleading
\emph{absence} of asks works by making a policy's own inferences point at the
wrong half-suit; a policy whose candidate set is already pruned to asks with
strictly positive hard-consistent probability gives that misdirection much less
to move. Robustness here was bought by removing a defect, not by adding a
model --- which is the more encouraging of the two possibilities, and also the
one with a lower ceiling.

The same panel records a loss. Against the feint --- which manufactures a
\emph{false} ask-legality certificate rather than withholding a true one ---
\vfive{} scores \numDecFeintFive\% against \vfour{}'s \numDecFeintFour\%, a
difference of $\numDecFeintDelta$ points that replicates at both seeds. It is
not an accident of tuning noise: the fit raised the weight on ``this player asked here''
from \numPriorThetaFour{} to \numPriorThetaFive{}, and the diagnosis had already
established that over-weighting the policy prior is exactly this exposure ---
while deleting the prior outright is worse still against the deceptive
archetypes themselves, by \numPriorDeleteCost{} points
[\numPriorDeleteCI]. Across the full twelve-style set \vfive{}'s worst case is the
\numDecWorstOpp{} at \numDecWorst\%, marginally below its mirror. No aggregate
figure in this report should be quoted without that table beside it.

\paragraph{What comes next, in order.} The ranking follows from the exposure
above rather than from the design document's build order.

\begin{enumerate}[leftmargin=1.4em,itemsep=1pt,topsep=2pt,parsep=0pt]
\item \textbf{M7, the online per-seat type posterior.} This is the principled fix
      for the feint, and \S\ref{sec:limitations} shows why nothing cheaper will
      do: a sweep of the prior weight over \numThetaSweepN{} values at two seeds
      does not identify the parameter. Shaped as a data-biased response
      \cite{johanson-dbr}, a manufactured certificate from a seat whose type
      posterior has drifted toward deceptive is discounted rather than believed,
      and the exploitation/robustness dial becomes explicit and reportable
      \cite{johanson-rnash} instead of being frozen at a fitted constant.
\item \textbf{M4, M5 and M3, in that order}, because each is the precondition of
      the next. A per-seat knowledge model is nearly free from a public
      transcript; it makes the target dimension scoreable, which is currently
      discarded outright --- the ask score opens by voiding its target argument,
      throwing away a mean of \numTargetBits{} free bits per ask, with at least
      two hard-indistinguishable targets on offer at \numTargetTwo\% of
      decisions and roughly \numTargetBitsGame{} bits per team per game in
      total; and only
      then can the net-information term of M3 be priced rather than gated, which
      \S\ref{sec:diag-information} shows is the only form in which it can ever
      fire.
\item \textbf{An exploitability probe.} The v0.4 study ran one; this one
      did not, and until it does \vfive{} claims no robustness property that a best
      response could contradict. Such a probe reports an achieved score that is
      a lower bound within its search class \cite{lisy-lbr}, and it is worth
      running both with and without M6, since a deterministic argmax sits at
      the deterministic leakage bound whatever its weights
      \cite{alvim-leakage}.
\item \textbf{A value model that is more than one feature.} Held-out $R^2$ is
      \numValueRSq{} for all sixteen features and \numValueRSqScore{} for a bias
      term plus the score differential alone, while a tree reaches
      \numValueRSqTree{}. Either rebuild it on rows collected at declaration
      points as well as ask points and refit it jointly with the policy vector,
      or retire it and score declarations directly on allocation probability.
      The measurement decides; the present state --- sixteen coefficients
      carrying one feature's worth of signal --- should not survive either way.
\end{enumerate}

\paragraph{The narrow claim, and the broader one.} The narrow claim of this
report is that a specific, reproducible failure of a specific agent was
diagnosed to an arithmetic cause in its ask score, that the cause was removed,
that the removal cost nothing measurable in aggregate win rate and left \vfour{}
marginally ahead on minimax regret, and that it helped substantially against one
class of deceptive opponent while hurting measurably against another.
The broader observation is
about evaluation. Every statistic in \S\ref{sec:diagnosis} was available to the
v0.4 study; none of it was collected, because the panel it was measured against
contained no opponent of its own strength, and the failure is a strong-opponent
phenomenon. The mirror belongs in the fitting panel, the pathology counters
belong in the commit gate, and the per-opponent breakdown with an explicit worst
case belongs in the abstract --- not because they are more rigorous, but because
an aggregate win rate against weaker opponents would have reported all of this
as a success.
